\documentclass[a4paper,11pt]{article}
\usepackage[utf8]{inputenc}
\usepackage[T1]{fontenc}
\usepackage[french]{babel}
\usepackage[left=2cm,right=2cm,top=2cm,bottom=2cm]{geometry}
\usepackage{xcolor}
\usepackage{hyperref}
\usepackage{fontawesome5}
\usepackage{parskip}

% --- Couleurs ---
\definecolor{primary}{RGB}{35, 55, 123}

\begin{document}

% --- EN-TÊTE ---
\begin{center}
    {\Large \textbf{Loïc Aron MBASSI EWOLO}}\\[4pt]
    \small
    \faMapMarker\ Cergy, France \quad
    \faPhone\ (+33) 7 59 52 33 98 \quad
    \faEnvelope\ \href{mailto:loicaron.mbassiewolo@ensea.fr}{loicaron.mbassiewolo@ensea.fr}
\end{center}
\vspace{0.5cm}

% --- DESTINATAIRE ---
\hfill
\begin{minipage}[t]{0.45\textwidth}
    \textbf{GE HealthCare}\\
    Service Recrutement\\
    283 rue de la Minière\\
    78350 Buc
\end{minipage}

\vspace{1cm}

\textbf{Objet :} Candidature Stage Développement Application Web de Revue d'Images Mammographiques

\vspace{0.5cm}

Madame, Monsieur,

Le développement web appliqué au domaine médical est un sujet qui me tient à cœur. Je travaille actuellement sur un projet de détection de cancer du poumon par classification d'images médicales, et j'ai contribué pendant deux ans à une plateforme de télémédecine. Votre offre de stage sur le développement d'une application de revue d'images mammographiques correspond à ce double intérêt. Actuellement en dernière année d'école d'ingénieur à l'\textbf{École Polytechnique de Yaoundé} et en double diplôme à l'\textbf{ENSEA}, je souhaite rejoindre votre équipe logicielle pour la santé de la femme.

Mon projet le plus pertinent pour votre offre est celui de \textbf{détection de cancer du poumon} que je mène à Polytechnique : classification d'images médicales (cas bénins/malins) et mise en place de pipelines de prétraitement pour améliorer la précision des modèles. Ce travail m'a familiarisé avec les problématiques de visualisation et d'analyse d'images médicales que vous rencontrez en mammographie.

J'ai également conçu un \textbf{Serious Game pour la formation médicale} sous Unity : un simulateur de consultation destiné aux internes en médecine. Ce projet m'a appris à travailler sur des interfaces utilisateur en collaboration avec des besoins métiers spécifiques, ce qui correspond à votre attente de collaboration avec l'équipe UI/UX.

Côté développement web, j'ai travaillé deux ans sur \textbf{SOSDOCTA}, une plateforme de télémédecine (réservation, paiement, dashboards médecins/patients), et j'ai développé une interface marketplace en \textbf{React.js} chez TOGETTECH. Mon expérience avec \textbf{TypeScript} sur un projet en stack MERN complète ce profil frontend.

Je suis disponible dès \textbf{avril 2026} pour un stage de 6 mois. Je serais ravi d'échanger sur vos besoins et de vous montrer mes réalisations lors d'un entretien.

Je vous prie d'agréer, Madame, Monsieur, l'expression de mes salutations distinguées.

\vspace{1cm}

\textbf{Loïc Aron MBASSI EWOLO}\\
\textit{Élève-Ingénieur 5\textsuperscript{ème} année - Polytechnique Yaoundé / ENSEA}

\end{document}
